\section{Results}
\label{sec:results}

We report on run \texttt{00000000c002daaaX} (experiment seed
$3{,}221{,}412{,}522$, \texttt{tarel} v0.0.10). Two cells are
analysed: $\texttt{binance:BTC-USDC}$ and $\texttt{kraken:BTC-USDC}$.
The third cell ingested under this run, $\texttt{coinbase:BTC-USDC}$,
is excluded for the reason given in Section~\ref{sec:data}.

\subsection*{Stage~1: bucket construction yield}
Table~\ref{tab:stage_one_yield} reports the per-(cell, clock) bucket
construction yield. $N_{\text{buckets}}$ is the count of windows the
bucketer enumerated, $N_{\text{dropped}}$ the count rejected by the
8-snapshot floor, $N_{\text{trunc}}$ the count that crossed a
60-minute parquet boundary and were truncated, and
$N_{\text{rec}}$ the count that reached the recurrence-plot writer
and produced an on-disk record. The yield column is
$N_{\text{rec}} / N_{\text{buckets}}$ as a percentage.

\input{content/sections-v2/04-results-discussion/tables/yield}

Two patterns are visible. First, the physical and event clocks behave
similarly across the two cells: yields are $\geq 99.9\%$ for binance
under both clocks and $\geq 85.8\%$ for kraken, with the writer
exercising thousands of buckets per (cell, clock) and the
minimum-snapshot floor pruning a single-digit fraction at most.
Second, the info\_time clock compresses the bucket population by one
to two orders of magnitude relative to the other two clocks: binance
has 50 info\_time buckets against 9{,}197 event buckets (a $184\times$
contraction); kraken's contraction is more moderate ($\sim 1.2\times$
event-to-info\_time) reflecting the venue's lower trade intensity.
The locked stage-2 batch size of 64 therefore sits inside binance's
info\_time bucket population only marginally, and the stage-2 grid
documented in Section~\ref{sec:methods} is partially sparse on
binance's info\_time and event clocks for that reason.

\subsection*{Stage~1: RQA characterisation}
Table~\ref{tab:stage_one_rqa} reports the Marwan-standard RQA scalars
(recurrence rate, determinism, laminarity, longest diagonal,
trapping time, entropy of the diagonal-line-length distribution)
aggregated as mean$\pm$std across the bucket population per (cell,
clock), at the locked target $RR^{*} = 5\%$. Buckets with $n < 16$
snapshots (the stage-2 batch viability floor) are excluded from
aggregation; the line-of-identity is excluded from the diagonal
counts following the Marwan convention. Diagonal and vertical line
minimum lengths are $l_{\min} = v_{\min} = 2$. The binance/info\_time
row is computed on the first $8$ of $40$ admitted buckets due to a
wallclock cap on the present analysis; the remaining $32$ buckets are
on disk and available for a full-cap re-run.

\input{content/sections-v2/04-results-discussion/tables/rqa}

The aggregate $\overline{RR}$ is on target ($\in [0.040, 0.050]$) for
all six (cell, clock) cells except kraken/event ($\overline{RR} =
0.026$), where the RR-bisection converges below target on the smaller
$n$ populations characteristic of this clock. $\overline{DET}$ and
$\overline{LAM}$ separate by clock more sharply than by cell: the
info\_time clock yields the highest determinism and laminarity at
both cells (binance: $\overline{DET} = 0.801$,
$\overline{LAM} = 0.896$; kraken: $0.615, 0.647$), reflecting the
volume-thresholded info\_time bucketer's tendency to span longer
quasi-stationary regimes. $\overline{L_{\max}}$ on
binance/info\_time is large ($90.1$) but proportional to the
much larger $n$ of those buckets ($n \sim 10^4$, vs. $n \sim 64$
for binance/event).

\subsection*{Stage~1: distributional structure}
Figure~\ref{fig:bucket_duration} shows the per-bucket duration
distribution across the six (cell, clock) panels.
Figure~\ref{fig:rqa_violins} shows the per-bucket distributions of
$DET$ and $LAM$ as violins faceted by clock, with one violin per cell
inside each panel. Figure~\ref{fig:mean_rp} shows the population-mean
recurrence plot per (cell, clock) at $RR^{*} = 5\%$, computed by
nearest-neighbour resampling each bucket's binary RP onto a common
$64 \times 64$ grid and averaging.

\begin{figure}[t]
  \centering
  \includegraphics[width=\linewidth]{content/sections-v2/04-results-discussion/figures/bucket_duration.png}
  \caption{Per-bucket duration distributions across the two cells
  and three clocks (log-scale $x$). The info\_time clock contracts the
  bucket population substantially relative to physical and event.}
  \label{fig:bucket_duration}
\end{figure}

\begin{figure}[t]
  \centering
  \includegraphics[width=\linewidth]{content/sections-v2/04-results-discussion/figures/rqa_violins.png}
  \caption{Per-bucket distributions of $DET$ (top row) and $LAM$
  (bottom row) at $RR^{*} = 5\%$, faceted by clock; one violin per
  cell within each panel. The info\_time clock concentrates mass at
  high $DET$ / $LAM$ for both cells.}
  \label{fig:rqa_violins}
\end{figure}

\begin{figure}[t]
  \centering
  \includegraphics[width=\linewidth]{content/sections-v2/04-results-discussion/figures/mean_rp_per_clock.png}
  \caption{Population-mean binary recurrence plot per (cell, clock),
  $64 \times 64$ resample, $RR^{*} = 5\%$.}
  \label{fig:mean_rp}
\end{figure}

\subsection*{Stage~2: training behaviour}
Table~\ref{tab:stage_two_summary} reports the per-(cell, clock)
stage-2 execution summary at the locked seed budget of $J{=}1$:
wallclock at the last logged checkpoint, generator-step count
($1{,}000$), critic-step count ($5{,}000$), final-100-step EM-gap
mean$\pm$std, GP-magnitude mean across all critic steps, and the
GP-magnitude maximum. Pairs without a checkpoint directory on disk
are marked ``--'' for the step counters and wallclock; pairs whose
\texttt{diagnostics.msgpack} was flushed but whose checkpoint was
cut (kraken/info\_time on this run) still expose EM-gap and GP
statistics from the streamed diagnostics.

\input{content/sections-v2/04-results-discussion/tables/stage_two_summary}

Figure~\ref{fig:em_gap} shows the EM-gap trajectory (raw and
rolling-100 smoothed) for each successfully trained
(cell, clock, seed) trio in this run. Four of the six
(cell, clock) pairs reached the first checkpoint at step $1{,}000$;
the two un-trained pairs (binance/event and binance/info\_time)
were not launched on the present hardware schedule and appear as
empty panels.

\begin{figure}[t]
  \centering
  \includegraphics[width=\linewidth]{content/sections-v2/04-results-discussion/figures/em_gap_trajectories.png}
  \caption{Stage-2 EM-gap convergence per trained (cell, clock); raw
  per-step values (light) and rolling-100 smooth (dark); dashed line
  marks the final-100-step mean reported in the inset.}
  \label{fig:em_gap}
\end{figure}

Figure~\ref{fig:loss_by_cell} aggregates the three primary
stage-2 signals (EM-gap, GP magnitude, and the critic-loss
approximation $\hat{L}_D \approx -\text{EM} + \lambda_{\mathrm{GP}}
\cdot \text{GP}$) per cell, overlaying the available clocks within
each panel. The view complements Figure~\ref{fig:em_gap} by
collapsing the (cell, clock) grid into per-cell trajectories and
exposes that kraken's three trained clocks converge to consistent
relative EM-gap orderings (event $<$ info\_time $<$ physical) and
that the GP magnitude across all four trained pairs remains
materially above the Lipschitz target of $1$ throughout the
$1{,}000$-step horizon.

\begin{figure}[t]
  \centering
  \includegraphics[width=\linewidth]{content/sections-v2/04-results-discussion/figures/loss_evolution_by_cell.png}
  \caption{Stage-2 loss-function evolution aggregated by cell: rows
  show EM-gap (rolling-100 smooth), gradient-penalty magnitude
  (symlog-y), and the critic-loss approximation
  $\hat{L}_D = -\text{EM} + \lambda_{\mathrm{GP}} \cdot \text{GP}$
  at the locked $\lambda_{\mathrm{GP}}{=}10$. Columns are cells;
  line style distinguishes the three clocks within each panel.}
  \label{fig:loss_by_cell}
\end{figure}

\paragraph{Partial Stability Score.}
The locked four-axis Stability Score
(\texttt{StabilityScoreSpec\_v0\_1\_0.json}) requires the
mode-coverage and critic-loss-curvature axes, which are populated
only at $\geq 5{,}000$ generator-step checkpoints. The present run's
training horizon stopped at the first checkpoint (step $1{,}000$),
so those two axes are empty for every trained pair and the locked
Score is not computable from this run. We report instead a
\emph{2-axis partial} Stability Score that uses the EM-gap and
GP-magnitude axes at equal weight ($0.5$ each), rank-normalised
across the four trained (cell, clock) pairs. Among the four trained
pairs, kraken/info\_time ranks highest ($0.83$), kraken/event next
($0.67$), kraken/physical next ($0.50$), and binance/physical lowest
($0.00$). The numeric values appear in
\texttt{figures/partial\_stability\_scores.json} and are not
report-comparable against article-1's locked 4-axis aggregate; the
ordering should be read as a within-run ranking only.
